\documentclass[11pt,a4paper]{article}
\usepackage[utf8]{inputenc}
\usepackage[T1]{fontenc}
\usepackage[spanish]{babel}
\usepackage[margin=2.5cm]{geometry}
\usepackage{booktabs}
\usepackage{array}
\usepackage{tabularx}
\usepackage{listings}
\usepackage{xcolor}
\usepackage{enumitem}
\usepackage{parskip}
\usepackage{fancyhdr}
\usepackage{amsmath}
\usepackage{amssymb}

\definecolor{codegray}{RGB}{245,245,245}
\definecolor{codeblue}{RGB}{0,60,180}
\definecolor{codegreen}{RGB}{0,120,0}

\lstdefinestyle{log}{
  basicstyle=\ttfamily\small,
  backgroundcolor=\color{codegray},
  frame=single,
  framerule=0.4pt,
  xleftmargin=8pt,
  numbers=none,
  breaklines=true
}

\pagestyle{fancy}
\fancyhf{}
\rhead{\textcolor{gray}{\small TD7: Ingeniería de Datos — UTDT}}
\lhead{\textcolor{gray}{\small Segundo Parcial Modelo 4}}
\rfoot{\thepage}
\renewcommand{\headrulewidth}{0.4pt}

\setlist[enumerate,1]{label=\textbf{\alph*)}}

\begin{document}

\begin{center}
  {\LARGE\bfseries TD7 — Ingeniería de Datos}\\[6pt]
  {\large Universidad Torcuato Di Tella}\\[4pt]
  {\Large\bfseries Segundo Parcial — Modelo 4}\\[4pt]
  {\normalsize Duración: 2 horas \quad|\quad Puntaje total: 100 puntos}
\end{center}
\noindent\rule{\linewidth}{0.8pt}
\vspace{2pt}

\noindent\textbf{Instrucciones:} Responda todos los ejercicios. Justifique cada respuesta; las afirmaciones sin justificación no suman puntaje. No está permitido el uso de material de consulta ni dispositivos electrónicos.

\vspace{6pt}
\noindent\rule{\linewidth}{0.4pt}

% ─────────────────────────────────────────────────────────────
\section*{Ejercicio 1 — Optimización en SQL Server \hfill\normalfont(30 puntos)}

Considere el siguiente esquema y consulta sobre una base en SQL Server:

\begin{center}
\textbf{Cliente}(\underline{\texttt{id\_cliente}}, \texttt{nombre}, \texttt{ciudad}) \qquad
\textbf{Venta}(\underline{\texttt{id\_venta}}, \texttt{id\_cliente}, \texttt{fecha}, \texttt{monto})
\end{center}

\begin{lstlisting}[style=log]
SELECT  c.ciudad, SUM(v.monto) AS total
FROM    Cliente c JOIN Venta v ON c.id_cliente = v.id_cliente
WHERE   c.ciudad = 'Rosario'
GROUP BY c.ciudad;
\end{lstlisting}

La tabla \textbf{Venta} tiene 10 millones de filas; \textbf{Cliente} tiene 50.000 filas. Existe un índice \textit{nonclustered} sobre \texttt{Cliente.ciudad} y el índice \textit{clustered} de cada tabla está sobre su clave primaria.

\begin{enumerate}
  \item (6 pts) En un plan de ejecución de SQL Server, ¿en qué dirección se lee el flujo de datos y qué representa el \textbf{porcentaje} mostrado en cada operador? Explique la diferencia entre un operador \textbf{Scan} y uno \textbf{Seek}, y cuándo cada uno es preferible.

  \vspace{4cm}

  \item (8 pts) Para acceder a las filas de \texttt{Cliente} con \texttt{ciudad = 'Rosario'}, ¿el optimizador usará un \textit{Index Seek} o un \textit{Clustered Index Scan}? ¿Y para leer \texttt{Venta}? Justifique en función de la selectividad de cada predicado y de los índices disponibles.

  \vspace{4.5cm}

  \item (9 pts) El optimizador debe elegir el algoritmo de \textbf{junta}. Complete y justifique: para cada algoritmo (\textbf{Nested Loops}, \textbf{Merge Join}, \textbf{Hash Join}) indique si requiere que las entradas estén ordenadas, si requiere un equi-join y en qué situación de tamaños relativos es óptimo. ¿Cuál esperaría que elija el optimizador para esta consulta y por qué?

  \vspace{6cm}

  \item (7 pts) La consulta tiene \texttt{SUM} con \texttt{GROUP BY}. Explique la diferencia entre los operadores \textbf{Stream Aggregate} y \textbf{Hash Aggregate}. ¿Cuál exige que los datos lleguen ordenados? Si creara un índice que provea ese orden, ¿qué operador de agregación esperaría ver y por qué?

  \vspace{4.5cm}
\end{enumerate}

\newpage

% ─────────────────────────────────────────────────────────────
\section*{Ejercicio 2 — Integración de Datos \hfill\normalfont(25 puntos)}

\begin{enumerate}
  \item (8 pts) Complete una tabla comparativa entre una base \textbf{transaccional (OLTP)} y un \textbf{Data Warehouse (OLAP)} en las dimensiones: optimización (escritura vs.\ lectura), normalización, qué información retiene (estado actual vs.\ historia) y perfil de latencia/throughput. Explique por qué \textbf{no} conviene correr reportes analíticos pesados directamente sobre la base transaccional de producción.

  \vspace{5.5cm}

  \item (9 pts) Describa la arquitectura \textbf{Medallion} (capas \textbf{Bronze}, \textbf{Silver} y \textbf{Gold}), indicando qué contiene cada capa. Mencione dos ventajas que aporta esta organización en capas (por ejemplo, frente a errores de datos).

  \vspace{5.5cm}

  \item (8 pts) Diseñe un \textbf{star schema} para la capa Gold de un \textit{e-commerce} que necesita reportar ventas por producto, cliente, fecha y canal. Indique la \textbf{fact table} (con sus claves y métricas) y las \textbf{dimension tables}. Explique la diferencia estructural entre fact y dimension tables (filas vs.\ columnas). Si la dirección de un cliente cambia, ¿qué significa \textbf{pisar} vs.\ \textbf{versionar} y qué se pierde con cada opción?

  \vspace{6cm}
\end{enumerate}

\newpage

% ─────────────────────────────────────────────────────────────
\section*{Ejercicio 3 — Arquitectura de Datos y Sistemas Distribuidos \hfill\normalfont(25 puntos)}

\begin{enumerate}
  \item (8 pts) Contraste las estrategias \textbf{Fully Incremental} y \textbf{Event Sourcing}: qué almacena cada una, cómo responde una consulta y cuál es el problema central de cada una (tolerancia a errores humanos en una; costo/escala de almacenamiento y cómputo en la otra).

  \vspace{5cm}

  \item (9 pts) Defina los niveles de consistencia: \textbf{fuerte}, \textbf{read-your-own-writes (RYOW)}, \textbf{de sesión} y \textbf{eventual}. Para un sistema donde un usuario actualiza su propio perfil y espera verlo actualizado al instante, pero a otros usuarios les puede llegar con retraso, ¿qué nivel alcanza y por qué?

  \vspace{6cm}

  \item (8 pts) Enuncie el \textbf{teorema CAP} (consistencia, disponibilidad, tolerancia a particiones). Explique por qué, ante una \textbf{partición de red}, un sistema distribuido debe elegir entre CP y AP. Dé un ejemplo de aplicación donde priorizaría C sobre A y otro donde priorizaría A sobre C.

  \vspace{5cm}
\end{enumerate}

\newpage

% ─────────────────────────────────────────────────────────────
\section*{Ejercicio 4 — Anonimización y Privacidad \hfill\normalfont(20 puntos)}

Un hospital quiere publicar el siguiente dataset para investigación. Los atributos \texttt{Edad} y \texttt{CP} (código postal) son \textbf{cuasi-identificadores}; \texttt{Enfermedad} es el atributo \textbf{sensible}.

\begin{center}
\begin{tabular}{cccl}
\toprule
\textbf{\#} & \textbf{Edad} & \textbf{CP} & \textbf{Enfermedad} \\
\midrule
1 & 29 & 2000 & Gripe \\
2 & 31 & 2001 & HIV \\
3 & 27 & 2002 & HIV \\
4 & 55 & 2010 & Diabetes \\
5 & 58 & 2011 & Gripe \\
6 & 52 & 2012 & Cáncer \\
\bottomrule
\end{tabular}
\end{center}

\begin{enumerate}
  \item (5 pts) Defina \textbf{PII}. Explique por qué eliminar solo el nombre y el DNI \textbf{no} garantiza el anonimato si quedan cuasi-identificadores como edad y código postal.

  \vspace{3.5cm}

  \item (8 pts) Explique la técnica de \textbf{$k$-anonimato}. Generalizando \texttt{Edad} (en rangos) y \texttt{CP} (truncando dígitos), proponga una transformación de la tabla que logre \textbf{$k=3$}. Muestre los grupos resultantes.

  \vspace{5cm}

  \item (7 pts) Explique qué es la \textbf{$\ell$-diversidad} y qué problema del $k$-anonimato resuelve. En la tabla original, ¿qué podría inferir un atacante sobre el grupo de los registros 2 y 3 aunque estuvieran $k$-anonimizados? ¿Cómo lo evitaría la $\ell$-diversidad?

  \vspace{4.5cm}
\end{enumerate}

\end{document}
